\section{Worker II: creative telescoping without boundary fiction}
\label{sec:telescoping}

\subsection{A different problem from polynomial antidifferencing}

Forge already searches for polynomial $q$ satisfying $q(n+1)-q(n)=f(n)$ for a
polynomial increment $f$. This is useful, but it is not a method for sums whose
summands contain factorial or binomial quotients. The present worker targets
\begin{equation}\label{eq:sum}
S_{p,q}(n)=\sum_{k=0}^{n}q^k\binom nk^p,
\qquad n\in\N,\quad p\in\{1,2,\ldots\},\quad q\in\Q.
\end{equation}
The concrete implementation bounds $p\le8$, recurrence order at most six, and
polynomial degrees through a specified search request. Its experimental corpus
uses $p\le4$. There is no claim that every possible request inside those decoder
bounds will succeed.

Creative telescoping seeks a recurrence for $S$ by proving a local identity for
the summand. It is a mature method rather than a new proposed theory. The
production recommendation is to reuse HolonomicFunctions for broader
multivariate telescoping and \code{ore_algebra} for operator computations,
while retaining a small independently checked certificate language~\cite{holonomic,ore}.
The prototype is intentionally more limited and easier to audit.

\subsection{Why a rational multiplier is not enough}

The usual presentation writes $G(n,k)=R(n,k)F(n,k)$ with a rational function $R$
and a hypergeometric term $F$. Away from poles, a rational identity can show
\[
\sum_i a_i(n)F(n+i,k)=G(n,k+1)-G(n,k).
\]
At a support boundary, however, $F$ may vanish while $R$ has a pole. Declaring
$R F=0$ there because ``the summand is zero'' is not a valid argument. Cancelling
a rational expression on an interior domain does not choose its extension at a
singular endpoint.

For example, even the relation
\[
\binom n{k-1}=\frac{k}{n-k+1}\binom nk
\]
requires care at $k=n+1$. Its left side is $1$, while the right presentation has
a zero denominator multiplied by a zero binomial coefficient. The shifted
binomial is the well-defined object. The rational presentation is only a tool
for algebra on the interior.

The safe certificate therefore stores the flux as a sum of \emph{shifted
summands}, not as an unevaluated singular product.

\subsection{The shifted-support ansatz}

Extend the binomial coefficient to a signed lower index by
\[
\binom nk=0\quad\text{when }k<0\text{ or }k>n,
\]
with its ordinary value otherwise. For an order $r\ge1$ recurrence, search for
univariate polynomials $a_0,\ldots,a_r\in\Q[n]$, with $a_r\ne0$, and bivariate
polynomials $P_0,\ldots,P_{r-1}\in\Q[n,k]$. Define, for $k\ge0$,
\begin{equation}\label{eq:flux}
G(n,k)=\sum_{j=0}^{r-1}P_j(n,k)q^k\binom{n+j}{k-1}^{p}.
\end{equation}
The desired summand identity is
\begin{equation}\label{eq:tel-id}
\sum_{i=0}^{r}a_i(n)q^k\binom{n+i}{k}^{p}
 =G(n,k+1)-G(n,k),\qquad 0\le k\le n+r.
\end{equation}
The two end fluxes are structurally zero:
\begin{equation}\label{eq:ends}
G(n,0)=0,\qquad G(n,n+r+1)=0.
\end{equation}
Indeed, the first uses the signed index $-1$, and in the second
$k-1=n+r>n+j$ for every $j<r$. No limiting argument is involved.

The fixed summation family is part of the authoritative problem. A certificate
contains only the coefficient polynomials. It cannot choose a different
summand, change $p$ or $q$, or assert its own support convention.

\subsection{Interior verification by a polynomial identity}

On $0\le k\le n$, every factor of
\[
D(n,k)=\prod_{h=1}^{r}(n+h-k)^p
\]
is positive. Put
\begin{align}
R_i(n,k)&=\prod_{h=1}^{i}(n+h)^p
          \prod_{h=i+1}^{r}(n+h-k)^p,\label{eq:Ri}\\
T_j(n,k)&=k^p\prod_{h=1}^{j}(n+h)^p
          \prod_{h=j+2}^{r}(n+h-k)^p.\label{eq:Tj}
\end{align}
Empty products are one. The binomial quotient identities give
\[
R_i=D\frac{\binom{n+i}k^p}{\binom nk^p},\qquad
T_j=D\frac{\binom{n+j}{k-1}^p}{\binom nk^p}.
\]
The checker verifies the exact sparse-polynomial equality
\begin{equation}\label{eq:interior}
\sum_{i=0}^{r}a_iR_i
-\sum_{j=0}^{r-1}qP_j(n,k+1)R_j
+\sum_{j=0}^{r-1}P_j(n,k)T_j=0.
\end{equation}
This implies \eqref{eq:tel-id} on the interior. Only nonzero binomial and
$D$ factors are used in the rational derivation. The factor $q^k$ is multiplied
back, \emph{never divided out}. Therefore the proof also covers $q=0$, with the
ordinary convention $0^0=1$.

All arithmetic in the acceptance checker is addition and multiplication of
sparse rational polynomials. Search-generated residual expressions are not
accepted on assertion: the checker reconstructs the residual from the
authoritative problem and the supplied coefficients, then requires its
canonical dictionary to be empty.

\subsection{Every exceptional boundary is a separate checked identity}

The interior covers $k\le n$, but \eqref{eq:tel-id} is needed through $n+r$.
For each $t=1,\ldots,r$, set $k=n+t$. Define the polynomial
\[
B_{j,t}(n)=
\begin{cases}
\displaystyle\binom{n+j}{j-t},&j-t\ge0,\\
0,&j-t<0.
\end{cases}
\]
When the lower index is nonnegative, it is fixed and the binomial is a
polynomial in $n$ with rational coefficients:
\[
\binom{n+j}{d}=\frac{1}{d!}\prod_{h=0}^{d-1}(n+j-h).
\]
The checker verifies, for every $t=1,\ldots,r$,
\begin{equation}\label{eq:boundary}
\begin{split}
0={}&\sum_{i=0}^{r}a_i(n)B_{i,t}(n)^p\\
&-q\sum_{j=0}^{r-1}P_j(n,n+t+1)B_{j,t}(n)^p\\
&+\sum_{j=0}^{r-1}P_j(n,n+t)B_{j,t-1}(n)^p.
\end{split}
\end{equation}
This is again a polynomial identity, now univariate. Multiplying it by
$q^{n+t}$ gives the exact boundary summand identity. There is no assumption
$q\ne0$. Complementing the binomial lower index is justified because $n\ge0$
and the nonnegative fixed lower indices are no larger than the upper indices.

\begin{theorem}[Boundary-safe telescoping]\label{thm:tel}
If $a_r\ne0$ and the polynomials satisfy \eqref{eq:interior} and all $r$
identities \eqref{eq:boundary}, then
\begin{equation}\label{eq:recurrence}
\sum_{i=0}^{r}a_i(n)S_{p,q}(n+i)=0\qquad(n\in\N).
\end{equation}
\end{theorem}
\begin{proof}
The interior calculation and the $r$ boundary calculations establish
\eqref{eq:tel-id} at every integer $k$ in $[0,n+r]$. Sum that identity over this
whole interval. Each summand $\binom{n+i}k^p$ vanishes above $n+i$, so its sum
is exactly $S_{p,q}(n+i)$, even though the different shifted sums have different
natural support endpoints. The flux telescopes to
$G(n,n+r+1)-G(n,0)=0$ by \eqref{eq:ends}. This proves \eqref{eq:recurrence}.
\end{proof}

The theorem is a deliberate alternative to a general-purpose rational
singularity solver. A restricted flux language makes the endpoints and support
coverage cheap to prove. Broader summands should enter through additional
proved support schemas, not by weakening this theorem's side conditions.

\subsection{How the search is implemented}

For chosen bounds $d_a,d_P$, write
\[
a_i(n)=\sum_{a=0}^{d_a}c_{i,a}n^a,\qquad
P_j(n,k)=\sum_{a+b\le d_P}d_{j,a,b}n^ak^b.
\]
Substitute these templates into \eqref{eq:interior} and every identity
\eqref{eq:boundary}. Equating coefficients produces a homogeneous rational
linear system. Its number of unknowns is
\begin{equation}\label{eq:ansatz-size}
(r+1)(d_a+1)+r\binom{d_P+2}{2}.
\end{equation}
The producer computes its exact nullspace using SymPy, selects a basis vector
with nonzero highest recurrence coefficient, clears rational denominators,
removes the common integer content, and chooses a positive leading coefficient
for $a_r$. It then sends the certificate through the independent checker.

Including the boundary equations in the search is important for completeness
within this ansatz. Computing only the interior nullspace and testing its
individual basis vectors against the boundaries could miss a useful linear
combination. The delivered implementation includes all those equations before
nullspace calculation. If every nullspace basis vector has $a_r=0$, so does
every linear combination, and this particular order and degree request has no
admissible solution. The worker reports \texttt{UNKNOWN\_FIXED\_ANSATZ}, not
``no recurrence exists'' or ``the theorem is false''.

For the fixed finite linear system, the selection of a top-nonzero nullspace
vector is complete. This does \emph{not} imply completeness of the union of
this particular flux language over all bounds, nor of creative telescoping
for arbitrary expressions. The number of monomials and coefficient bit sizes
can also make an exact nullspace calculation expensive. The quartic-power
example below already takes most of the measured telescoping search time.

\subsection{Recovered certificates, not just guessed recurrences}

\paragraph{Binomial theorem.}
For $p=1$, one certificate is
\[
a_0=-(1+q),\qquad a_1=1,\qquad P_0=-1.
\]
Thus $S(n+1)=(1+q)S(n)$ and $S(0)=1$. The flux is
$-q^k\binom n{k-1}$. This is a simple sanity case, not evidence that Forge
needs to rediscover an identity already available in Mathlib's binomial
library~\cite{mathlib-choose}. A production router should prefer a direct
library theorem when it matches.

\paragraph{Squared binomials.}
For $p=2,q=1$, the search returns
\begin{equation}\label{eq:squarecert}
a_0=-2(2n+1),\quad a_1=n+1,\quad P_0=2k-3n-3.
\end{equation}
The recurrence is
$(n+1)S(n+1)=2(2n+1)S(n)$. Combined with the elementary factorial recurrence
of $\binom{2n}n$ and the initial value $1$, it proves
\[
\sum_{k=0}^{n}\binom nk^2=\binom{2n}n.
\]
The artifact establishes the summation certificate, not an already integrated
Lean proof of this closed form. At $n=0$, the boundary cell $k=1$ contributes
exactly the term needed for the recurrence. Dropping it is not harmless
simplification of a zero tail.

More generally, the following formula can be checked by the displayed
polynomial identities for every rational $q$:
\begin{align*}
a_0&=(q-1)^2(n+1),&a_1&=-(q+1)(2n+3),&a_2&=n+2,\\
P_0&=(q-1)(n+1),&P_1&=2k-3n-6.
\end{align*}
Substitution into \eqref{eq:interior} and \eqref{eq:boundary} proves the
parametric statement by polynomial expansion. The executed search fixes $q$
before solving its rational system; it did not infer this as a universal
parameterised theorem. A separate supplied SymPy script also expands all three
residuals with symbolic $q$ and obtains zero; this check is outside the main
rational-certificate corpus. The corpus checks nine order-two rational instances,
including $0$, negative values, and nonintegral values, in addition to the
order-one $q=1$ instance.

\paragraph{Alternating squared binomials.}
For $q=-1$ the order-two recurrence becomes
\[
(n+2)S(n+2)+4(n+1)S(n)=0,
\]
with flux polynomials
\[
P_0=-2(n+1),\qquad P_1=2k-3n-6.
\]
The initial values are $S(0)=1,S(1)=0$. Consequently odd values vanish, and the
even subsequence is $S(2m)=(-1)^m\binom{2m}m$. The stride-two nature of this
example illustrates why a worker should retain the whole recurrence operator,
including a zero middle coefficient, rather than assume every identity reduces
to a first-order ratio.

\paragraph{Franel numbers.}
For $p=3,q=1$, the emitted operator is
\begin{equation}\label{eq:franel}
(n+2)^2S(n+2)-(7n^2+21n+16)S(n+1)-8(n+1)^2S(n)=0.
\end{equation}
Its explicit safe flux uses
\begin{align*}
P_0(n,k)&=-4(n+1)^2,\\
P_1(n,k)&=-6k^2+15kn+30k-10n^2-40n-40.
\end{align*}
This request has $r=2,d_a=d_P=2$, a $54\times21$ search matrix, nine boundary
rows, and a one-dimensional nullspace. A smaller first-order request with
$d_a=d_P=1$ has nullity zero and is preserved as an unknown result.
The initial values are $1,2$, and the leading coefficient $(n+2)^2$ has no
natural zero. Thus this particular recurrence needs no additional singular
seed values.

\paragraph{A larger, nonminimal certificate.}
For $p=4,q=1$, the chosen bounds $r=d_a=d_P=3$ produce a valid order-three
operator. Its last coefficient is $7(n+3)^3$; the full operator and three flux
polynomials are stored in the JSON corpus. The $159\times46$ matrix has a
two-dimensional nullspace. This is an acceptance of a correct recurrence, not
a claim that order three is minimal. Requiring minimal order in the checker
would add cost without strengthening the theorem established by a valid
annihilator.

\subsection{Lean-specific traps and the initial support library}

The most important translation trap is the lower index $k-1$. In Lean's
\code{Nat}, subtraction is truncated: $0-1=0$. Therefore translating
$\binom n{k-1}$ with a natural lower index would make the supposedly zero
left flux nonzero. The safe interface needs a signed lower index, for example
an explicitly defined \code{chooseZ : Nat -> Int -> Rat}, zero for negative
indices, with proved connections to \code{Nat.choose} on nonnegative inputs.
Alternatively, pattern match on $k$ before using natural subtraction. Either
choice must be reflected in the denotation proof.

An initial Lean support library should prove: zero outside the binomial
support; the interior shift-ratio identities with their inequalities; the
complement symmetry needed at the top boundary; the fixed-lower-index
polynomial product formula; finite-sum interchange; and the telescoping
endpoint theorem. Mathlib's choose and finite-sum libraries supply relevant
building blocks, but the exact certificate theorem still has to be written and
compiled. A generic \code{field_simp} invocation is not a substitute for the
missing domain proofs.

The source reifier must also distinguish the finite range
$\{0,\ldots,n\}$ from $\{0,\ldots,n-1\}$. In a \code{Finset.range} encoding,
\eqref{eq:sum} uses \code{n + 1}. The common support used in the proof has
length $n+r+1$. Off-by-one changes to either range alter the theorem.

The first support schema is intentionally finite. Infinite sums would require
convergence and limiting boundary proofs. Products of factorials with several
moving bounds require a richer support partition. Neither is admitted by the
present checker; neither should be approximated by checking a large finite
range and silently invoking the finite theorem.
